\documentclass{article}
\usepackage{geometry}
\geometry{margin=1.5cm, vmargin={0pt,1cm}}
\setlength{\topmargin}{-1cm}
\setlength{\headheight}{12.5pt}
\setlength{\paperheight}{29.7cm}
\setlength{\textheight}{25.3cm}

% useful packages.
% \usepackage[UTF8]{ctex}
\usepackage{fancyhdr}
\usepackage{layout}

% import common macros
% % % % % % % % % % % % % % % % % % % % % % % % % % % % % % % %
% Common Macros for LaTeX
% % % % % % % % % % % % % % % % % % % % % % % % % % % % % % % %

% Useful packages
\usepackage{amsfonts}
\usepackage{amsmath}
\usepackage{amssymb}
\usepackage{amsthm}
\usepackage{enumerate}
\usepackage{graphicx}
\usepackage{multicol}
\usepackage[ruled,linesnumbered]{algorithm2e}

% Math operators and common symbols
\newcommand{\dif}{\mathrm{d}}
\newcommand{\innerProd}[1]{\left\langle #1 \right\rangle}
\newcommand{\difFrac}[2]{\frac{\dif #1}{\dif #2}}
\newcommand{\pdfFrac}[2]{\frac{\partial #1}{\partial #2}}
\newcommand{\T}{\mathrm{T}}
\newcommand{\norm}[1]{\left\| #1 \right\|}
\newcommand{\abs}[1]{\left| #1 \right|}

% Vector and Matrix shortcuts
\newcommand{\vecTwo}[2]{
  \begin{bmatrix}
    #1 \\
    #2
\end{bmatrix}}
\newcommand{\vecThree}[3]{
  \begin{bmatrix}
    #1 \\
    #2 \\
    #3
\end{bmatrix}}
\newcommand{\vecTwoH}[2]{
  \begin{bmatrix}
    #1 & #2
  \end{bmatrix}
}
\newcommand{\vecThreeH}[3]{
  \begin{bmatrix}
    #1 & #2 & #3
  \end{bmatrix}
}

% Common fields
\newcommand{\R}{\mathbb{R}}
\newcommand{\C}{\mathbb{C}}
\newcommand{\Z}{\mathbb{Z}}
\newcommand{\N}{\mathbb{N}}

% Solutions and proofs
\newenvironment{solution}{
\begin{proof}[解]}{
\end{proof}}

% Utility to get environment variables (requires lualatex)
\newcommand{\getEnv}[2]{\directlua{tex.print(os.getenv("#1") or "#2")}}


% Name and uID
\newcommand{\name}{\getEnv{NAME}{NAME}}
\newcommand{\uid}{\getEnv{UID}{UID}}
\newcommand{\course}{Point Set Topology}
\newcommand{\hwNum}{5}

\begin{document}

\pagestyle{fancy}
\fancyhead{}
\lhead{\zhstr{\name} (\uid)}
\chead{\course\ \#\hwNum}
\rhead{\today}

\section{Question 1}

Let $(X,d)$ be a metric space and define
\[
  d^*(x,y)=\min\{d(x,y),1\}.
\]
Show that: (i) $d^*$ is a metric; (ii) the induced topology $\mathcal{T}^*$
equals the topology $\mathcal{T}$ induced by $d$.
Also decide whether there always exists $A>0$ such that
\[
  A\,d(x,y)\le d^*(x,y)\quad (\forall x,y\in X).
\]

\begin{solution}
  \textbf{(i) $d^*$ is a metric.}

  Non-negativity and symmetry are immediate from the definition of
  $\min$ and from the same properties of $d$.
  Also,
  \[
    d^*(x,y)=0 \iff \min\{d(x,y),1\}=0 \iff d(x,y)=0 \iff x=y.
  \]
  So only the triangle inequality needs proof.
  For $x,y,z\in X$,
  \[
    d^*(x,z)=\min\{d(x,z),1\}\le \min\{d(x,y)+d(y,z),1\}.
  \]
  Hence it is enough to show
  \[
    \min\{a+b,1\}\le \min\{a,1\}+\min\{b,1\}\quad (a,b\ge 0).
  \]
  If $a\ge 1$ or $b\ge 1$, the right-hand side is at least $1$, while
  the left-hand side is at most $1$.
  If $a<1$ and $b<1$, then the right-hand side is $a+b$, so
  \[
    \min\{a+b,1\}\le a+b=\min\{a,1\}+\min\{b,1\}.
  \]
  Therefore $d^*$ satisfies the triangle inequality.

  \textbf{(ii) $\mathcal{T}^*=\mathcal{T}$.}

  If $0<r\le 1$, then
  \[
    B_{d^*}(x,r)=\{y: d^*(x,y)<r\}=\{y:d(x,y)<r\}=B_d(x,r).
  \]
  Thus small open balls are identical for both metrics.

  Now prove both inclusions.

  $\mathcal{T}\subseteq \mathcal{T}^*$: let $U\in \mathcal{T}$ and $x\in U$. Choose $\varepsilon>0$
  with $B_d(x,\varepsilon)\subseteq U$.
  Set $r=\min\{\varepsilon,1/2\}$.
  Then $0<r\le 1$ and
  \[
    B_{d^*}(x,r)=B_d(x,r)\subseteq B_d(x,\varepsilon)\subseteq U.
  \]
  So $U\in \mathcal{T}^*$.

  $\mathcal{T}^*\subseteq \mathcal{T}$: let $U\in \mathcal{T}^*$ and $x\in U$. Choose $r>0$ with
  $B_{d^*}(x,r)\subseteq U$.
  Set $s=\min\{r,1/2\}>0$.
  Then $B_d(x,s)=B_{d^*}(x,s)\subseteq B_{d^*}(x,r)\subseteq U$, hence $U\in \mathcal{T}$.

  So $\mathcal{T}^*=\mathcal{T}$.

  \textbf{(iii) Existence of a uniform $A>0$.}

  Not always.
  Take $X=\R$ with the usual metric $d(x,y)=|x-y|$. If such $A>0$
  existed, then for all $x,y$,
  \[
    A|x-y|\le d^*(x,y)=\min\{|x-y|,1\}\le 1.
  \]
  Choose $|x-y|>1/A$. Then $A|x-y|>1$, contradiction.
  Hence in general no such $A$ exists.
\end{solution}

\section{Question 2}

Let
\[
  \R^\omega=\prod_{i\in\N}\R,
\]
and let $\R^\omega_{\mathrm{ez}}$ be the set of eventually-zero sequences.
Find interior and closure of $\R^\omega_{\mathrm{ez}}$ in $\R^\omega$ under:
(1) product topology $\mathcal{T}_p$;
(2) box topology $\mathcal{T}_b$.

\begin{solution}
  \textbf{A. In $(\R^\omega,\mathcal{T}_p)$:}

  A basic neighborhood has the form
  \[
    \prod_{i\in\N}U_i,
  \]
  where $U_i\subseteq\R$ is open and $U_i=\R$ for all but finitely many $i$.

  \emph{Interior is empty.}
  Let $x\in \R^\omega_{\mathrm{ez}}$ and let $N=\prod_i U_i$ be any
  basic neighborhood of $x$.
  Only finitely many coordinates are restricted. Let $F$ be that finite set.
  Define
  \[
    y_i=
    \begin{cases}
      x_i,& i\in F,\\
      1,& i\notin F.
    \end{cases}
  \]
  Then $y\in N$ (because unrestricted coordinates allow any real
  value), but $y$ has infinitely many nonzero entries, so $y\notin
  \R^\omega_{\mathrm{ez}}$.
  Hence no neighborhood of $x$ is contained in
  $\R^\omega_{\mathrm{ez}}$; therefore
  \[
    \operatorname{int}_{\mathcal{T}_p}(\R^\omega_{\mathrm{ez}})=\varnothing.
  \]

  \emph{Closure is all of $\R^\omega$.}
  Take any $x\in\R^\omega$ and any basic neighborhood $N=\prod_i U_i$ of $x$.
  Again, only finitely many coordinates are restricted; call that set $F$.
  Define
  \[
    z_i=
    \begin{cases}
      x_i,& i\in F,\\
      0,& i\notin F.
    \end{cases}
  \]
  Then $z\in N$ and $z$ is eventually zero.
  So every nonempty basic neighborhood meets $\R^\omega_{\mathrm{ez}}$.
  Hence
  \[
    \overline{\R^\omega_{\mathrm{ez}}}^{\,\mathcal{T}_p}=\R^\omega.
  \]

  \textbf{B. In $(\R^\omega,\mathcal{T}_b)$:}

  A basic neighborhood is still $\prod_i U_i$, but now every
  coordinate may be restricted.

  \emph{Interior is empty.}
  Fix $x\in\R^\omega_{\mathrm{ez}}$ and a box-basic neighborhood
  $N=\prod_i U_i$ of $x$.
  For all sufficiently large $i$, $x_i=0$, and since each $U_i$ is
  open containing $0$, choose $t_i\in U_i$ with $t_i\ne 0$.
  Define
  \[
    y_i=
    \begin{cases}
      x_i,& x_i\ne 0,\\
      t_i,& x_i=0.
    \end{cases}
  \]
  Then $y\in N$, and $y_i\ne 0$ for infinitely many $i$, so $y\notin
  \R^\omega_{\mathrm{ez}}$.
  Thus
  \[
    \operatorname{int}_{\mathcal{T}_b}(\R^\omega_{\mathrm{ez}})=\varnothing.
  \]

  \emph{Closure equals $\R^\omega_{\mathrm{ez}}$.}
  It suffices to show the complement is open.
  Take $x\notin \R^\omega_{\mathrm{ez}}$; then
  \[
    I=\{i\in\N: x_i\ne 0\}
  \]
  is infinite.
  For each $i\in I$, choose an open interval $V_i\ni x_i$ with $0\notin V_i$.
  For $i\notin I$, set $V_i=\R$.
  Then $N=\prod_i V_i$ is a box-basic neighborhood of $x$.
  Any $y\in N$ satisfies $y_i\ne 0$ for all $i\in I$, so $y$ has
  infinitely many nonzero coordinates.
  Hence $y\notin \R^\omega_{\mathrm{ez}}$.
  Therefore $N\subseteq \R^\omega\setminus \R^\omega_{\mathrm{ez}}$,
  and the complement is open.
  So $\R^\omega_{\mathrm{ez}}$ is closed in $\mathcal{T}_b$, and
  \[
    \overline{\R^\omega_{\mathrm{ez}}}^{\,\mathcal{T}_b}=\R^\omega_{\mathrm{ez}}.
  \]
\end{solution}

\section{Question 3}

Is $(\R^\omega,\mathcal{T}_p)$ metrizable? Is it connected?

\begin{solution}
  \textbf{Metrizable: yes.}
  Define
  \[
    d_\omega(x,y)=\sup_{i\in\N}\frac{d^*(x_i,y_i)}{i+1},
  \]
  where $d^*(u,v)=\min\{|u-v|,1\}$ on $\R$.
  Since $0\le d^*(x_i,y_i)\le 1$, the supremum is finite and $d_\omega\le 1$.

  Metric axioms follow from coordinatewise properties of $d^*$.
  For triangle inequality,
  \[
    \frac{d^*(x_i,z_i)}{i+1}
    \le \frac{d^*(x_i,y_i)+d^*(y_i,z_i)}{i+1},
  \]
  then take supremum over $i$:
  \[
    d_\omega(x,z)\le d_\omega(x,y)+d_\omega(y,z).
  \]
  So $d_\omega$ is a metric.

  Now compare topologies.

  Every $d_\omega$-ball is $\mathcal{T}_p$-open:
  fix $x$ and $\varepsilon>0$. Choose $N$ with $1/(N+1)<\varepsilon/2$.
  If for $1\le i\le N$ we require
  \[
    |y_i-x_i|<\frac{\varepsilon(i+1)}{2},
  \]
  then for $i\le N$,
  \[
    \frac{d^*(x_i,y_i)}{i+1}\le \frac{|x_i-y_i|}{i+1}<\frac\varepsilon2,
  \]
  and for $i>N$,
  \[
    \frac{d^*(x_i,y_i)}{i+1}\le \frac{1}{i+1}<\frac\varepsilon2.
  \]
  Hence $d_\omega(x,y)<\varepsilon$. This gives a cylinder
  neighborhood inside the ball.

  Every $\mathcal{T}_p$-basic neighborhood contains a $d_\omega$-ball:
  let
  \[
    N=\left(\prod_{i\in F}U_i\right)\times\left(\prod_{i\notin F}\R\right)
  \]
  be basic with finite $F$, and $x\in N$.
  Choose $\delta_i\in(0,1)$ such that
  $(x_i-\delta_i,x_i+\delta_i)\subseteq U_i$ for $i\in F$.
  Set
  \[
    \delta=\min_{i\in F}\frac{\delta_i}{i+1}>0.
  \]
  If $d_\omega(x,y)<\delta$, then for each $i\in F$,
  \[
    d^*(x_i,y_i)<(i+1)\delta\le \delta_i<1,
  \]
  so $|x_i-y_i|=d^*(x_i,y_i)<\delta_i$, hence $y_i\in U_i$.
  Thus $y\in N$.
  So topologies coincide: $\mathcal{T}_{d_\omega}=\mathcal{T}_p$.

  \textbf{Connected: yes (in fact path connected).}
  For $x=(x_i)$ and $y=(y_i)$, define
  \[
    H:[0,1]\to\R^\omega,\qquad H(t)=((1-t)x_i+ty_i)_i.
  \]
  Each coordinate map $t\mapsto (1-t)x_i+ty_i$ is continuous in $\R$.
  Therefore $H$ is continuous for the product topology (continuity of
  all coordinate projections).
  Also $H(0)=x$, $H(1)=y$.
  So any two points are joined by a path; hence $(\R^\omega,\mathcal{T}_p)$ is connected.
\end{solution}

\section{Question 4}

Let $(\R^\omega,\mathcal{T}_b)$ and let
\[
  U=\{x=(x_i): \exists C_x>0\text{ such that }|x_i|\le C_x\ \forall i\}
\]
(the set of bounded sequences).
Prove that $U$ and $\R^\omega\setminus U$ form a nontrivial
separation, and conclude $(\R^\omega,\mathcal{T}_b)$ is disconnected.

\begin{solution}
  First, both sets are nonempty:
  \[
    (0,0,\dots)\in U,
  \]
  while
  \[
    (1,2,3,\dots)\notin U.
  \]
  Clearly
  \[
    U\cap(\R^\omega\setminus U)=\varnothing,
    \qquad
    U\cup(\R^\omega\setminus U)=\R^\omega.
  \]
  So it remains to show both are open in $\mathcal{T}_b$.

  \textbf{$U$ is open.}
  Take $x\in U$. Choose $C_x$ with $|x_i|\le C_x$ for all $i$.
  Define box neighborhood
  \[
    N_x=\prod_{i\in\N}(x_i-1,x_i+1).
  \]
  If $y\in N_x$, then
  \[
    |y_i|\le |x_i|+1\le C_x+1\quad(\forall i),
  \]
  so $y$ is bounded, i.e. $y\in U$.
  Hence $N_x\subseteq U$, thus $U$ is open.

  \textbf{$\R^\omega\setminus U$ is open.}
  Take $x\notin U$ (so $x$ is unbounded).
  Use the same box neighborhood
  \[
    N_x=\prod_{i\in\N}(x_i-1,x_i+1).
  \]
  For any $y\in N_x$ and each $i$,
  \[
    |y_i|\ge |x_i|-1.
  \]
  Given $M>0$, unboundedness of $x$ gives some $i$ with $|x_i|>M+1$.
  Then $|y_i|>|x_i|-1>M$.
  So $y$ is unbounded, hence $y\notin U$.
  Therefore $N_x\subseteq \R^\omega\setminus U$, and the complement is open.

  Thus $U$ and $\R^\omega\setminus U$ are two disjoint nonempty open
  sets whose union is the whole space: a nontrivial separation.
  So $(\R^\omega,\mathcal{T}_b)$ is disconnected.
\end{solution}

\section{Question 5}

Let
\[
  X=\{(x,0):x\in\R\}\cup\{(x,1/x):x>0\}\subseteq\R^2
\]
with the subspace topology from $\R^2$. Determine whether $X$ is connected.

\begin{solution}
  Define
  \[
    f:X\to\R,\qquad f(x,y)=xy.
  \]
  This is continuous as the restriction to $X$ of the continuous
  multiplication map on $\R^2$.

  Now compute values on the two pieces of $X$:
  \[
    (x,0)\in X \implies f(x,0)=0,
  \]
  \[
    (x,1/x)\in X,\ x>0 \implies f(x,1/x)=1.
  \]
  Hence
  \[
    f(X)=\{0,1\}.
  \]
  Suppose $X$ were connected. Then its continuous image $f(X)$ would
  have to be connected in $\R$.
  But $\{0,1\}$ is not connected.
  Contradiction.

  Therefore $X$ is disconnected.
\end{solution}

\end{document}
